\section{Continuum Limit and Emergent Metric}
  \label{sec:continuum-limit-and-emergent-metric}

  The spectral distance provides a relational notion of proximity.
  We now identify the conditions under which it admits a smooth continuum
  approximation and an effective metric description.

  Consider families of relational graphs whose Laplacian spectra exhibit
  sufficient regularity.
  No manifold or dimensionality is assumed.
  The continuum regime is defined operationally as one in which spectral
  distances vary smoothly and admit consistent local approximations.

  A central requirement is local factorizability:
  the relational system decomposes approximately into weakly coupled
  subsystems, enabling local definitions of spectral distance.
  This property underlies effective locality and permits the introduction
  of auxiliary coordinate charts.

  Continuum reconstruction is non-unique.
  Spectral geometry is generically non-injective:
  distinct relational configurations may induce identical effective metrics,
  and a given configuration may admit multiple equivalent embeddings.
  This ambiguity is structural.

  When factorizability and spectral regularity hold,
  the spectral distance admits the local quadratic form
  \begin{equation}
    d_{\mathrm{spec}}(i,j)^2
    \;\approx\;
    g_{\mu\nu}(x)\,\Delta x^\mu \Delta x^\nu ,
  \end{equation}
  where $g_{\mu\nu}(x)$ summarizes the leading-order behavior of
  spectral distances.
  The metric tensor is derived, not fundamental.

  The approximation is valid only when higher-order spectral corrections
  remain subdominant.
  Outside this regime, no smooth geometric interpretation is assumed.

  The framework remains purely kinematic.
  No assumptions on dynamics, causality, or temporal ordering are required
  for metric emergence.

  \input{04-continuum-limit-and-emergent-metric/001-spectral-admissibility}

  Admissibility conditions may involve monotonic spectral filters,
  reflecting the insensitivity of continuum structure to
  fine-grained spectral rearrangements beyond a given resolution.

  The effective metric arises solely from relational spectral data
  and requires no additional fundamental fields.
